% Chapter 16, generated by make_chapters.py from papers/Ch16_Bio_14_Snarl_Schur/anccft27.tex.
% Source paper: Bio 14
\chapter[The BEST determinant along the snarl tree]{The BEST determinant along the snarl tree: Schur complements, boundary forests, and linear-time counting}\label{chap:ch16}
\noindent{\small\emph{Paper Bio~14 of the series. Verification script and transcript: \texttt{papers/Ch16\_Bio\_14\_Snarl\_Schur/}.}}\par\medskip
\begin{chaptersummary}
The number of sequences that spell a given single-stranded genome graph is a determinant, by the BEST
theorem: $\ec(G)=t_r(G)\prod_v(d^+_v-1)!$, where $t_r(G)=\det\Lt$ is a reduced Laplacian. Genome and
variation graphs are locally tangled but globally chain-like, organized by their snarl tree. A proposal
for this paper was to count Eulerian circuits by contracting a ``snarl tensor network'' along that tree.
What survives is simpler and exact. Eliminating the interior $I$ of a snarl is taking a Schur complement:
$\det\Lt=\det\Lt_{II}\cdot\det(\Lt/\Lt_{II})$. The Schur complement is again a reduced Laplacian, of a
weighted digraph on the snarl's boundary. And $\det\Lt_{II}$ counts the spanning forests of the interior
rooted into the boundary. The ``boundary tensor'' of a snarl is thus an $|B|\times|B|$ Laplacian and
nothing more, because the determinant is multiplicative. Eliminating snarls in the order of the tree
keeps the front bounded and the work linear in the size of the graph. On synthetic genome graphs of up to
$39\,714$ vertices the front stays at $8$, and the time per vertex stays at $9$--$11\,\mu$s. The
snarl-order determinant agrees with dense Bareiss elimination, with forced-arc contraction, and with
minimum-degree elimination, the last modulo two primes. On the de Bruijn graphs of $\phi$X174 the front
is $196$ at $k=7$ and $2$ at $k=12$. None of this helps the double-stranded count, which is
$\sP$-hard (Chapter~\ref{chap:ch10}). The proposal's claims about polyploid phasing and specific human loci are not
addressed.
\end{chaptersummary}

\medskip
\noindent\textbf{Keywords.} BEST theorem, matrix-tree theorem, Schur complement, Kron reduction,
snarl decomposition, genome graphs, variation graphs, Gaussian elimination, fill-in, genome assembly.

\section{The question}\label{ch16:sec:q}

A single-stranded genome graph $G$ is a balanced multidigraph. Its vertices are the $k$-mers and its arcs
the $(k+1)$-mers, with multiplicity. The circular sequences it spells are its Eulerian circuits, and the
BEST theorem counts them \cite{BEST}:
\[
\ec(G)=t_r(G)\prod_v(d^+_v-1)!,\qquad t_r(G)=\det\Lt,
\]
where $\Lt$ is the out-Laplacian with the row and column of a root $r$ deleted. For one strand, then,
counting is already polynomial. The issue is size: a pangenome graph has millions of vertices, and a
dense determinant is cubic.

Genome graphs are not dense. Variation graphs decompose into \emph{snarls}, subgraphs separated from the
rest by a pair of boundary nodes, nested in a snarl tree \cite{Pat18}. The proposal for this paper was to
turn each snarl into a ``local transition tensor'' and contract a tree tensor network, with bond
dimensions truncated by singular value decomposition. We keep the geometry and replace the machinery. The
right object is the Schur complement, it needs no truncation, and it is exact.

\section{Schur complements of Laplacians}\label{ch16:sec:schur}

Split the non-root vertices into an interior $I$ and the rest $B$, and write $\Lt$ in blocks.

\begin{proposition}[elimination]\label{ch16:prop:schur}
If $\Lt_{II}$ is invertible, then
\[
\det\Lt=\det\Lt_{II}\cdot\det S,\qquad S=\Lt/\Lt_{II}=\Lt_{BB}-\Lt_{BI}\Lt_{II}^{-1}\Lt_{IB}.
\]
$S$ is again a reduced Laplacian: its off-diagonal entries are $\le0$ and its row sums are $\ge0$, the
deficit being the weight of arcs to the root.
\end{proposition}

\begin{proof}
The determinant formula is the block factorization. $\Lt$ is a nonsingular M-matrix with non-negative
row sums. Both properties pass to Schur complements, which is the Kron reduction of network
theory \cite{DB13}. Concretely, row $b$ of $S$ is row $b$ of $\Lt$ with each arc into $I$ replaced by
the distribution of its first exit from $I$. That distribution is non-negative and has total mass at
most $1$.
\end{proof}

\begin{proposition}[boundary forests]\label{ch16:prop:forest}
$\det\Lt_{II}$ is the number of spanning forests of $I$ in which every vertex of $I$ chooses one out-arc
and every tree is rooted outside $I$, i.e.\ in $B\cup\{r\}$.
\end{proposition}

\begin{proof}
This is the all-minors matrix-tree theorem for the principal minor on $I$ \cite{Cha82}. Choosing one
out-arc per vertex of $I$ gives a functional graph, and the choice is counted iff it has no cycle
inside $I$.
\end{proof}

\begin{lemma}[pivots]\label{ch16:lem:piv}
If $G$ is strongly connected, every principal minor of $\Lt$ is positive. Hence Gaussian elimination of
$\Lt$ in any order of vertices has positive pivots, and $\det\Lt$ is their product.
\end{lemma}

\begin{proof}
By Proposition~\ref{ch16:prop:forest}, the minor on a set $J$ counts forests of $J$ rooted outside $J$. Strong
connectivity gives every vertex of $J$ a path out of $J$, hence at least one such forest.
\end{proof}

So a snarl with boundary $B$ contributes a scalar $\det\Lt_{II}$, the number of its boundary-rooted
forests, and hands its parent an $|B|\times|B|$ Laplacian. That Laplacian is the snarl's ``boundary
tensor''. It is a matrix, not a higher tensor, because the determinant is multiplicative. Its size is the
boundary size, and there is nothing to truncate.

\section{Elimination along the snarl tree}\label{ch16:sec:tree}

Order the vertices so that every snarl's interior comes before its boundary, children before parents:
a post-order of the snarl tree. Eliminating a vertex $v$ costs one division and $|\mathrm{row}(v)|\cdot
|\mathrm{col}(v)|$ updates, where the row and column are those of the current matrix. The \emph{front} at
$v$ is $w_v=|\mathrm{row}(v)|+|\mathrm{col}(v)|$.

\begin{theorem}\label{ch16:thm:linear}
Elimination in snarl-tree order computes $t_r(G)$ exactly with $O\bigl(\sum_vw_v^2\bigr)$ field
operations. If every snarl has at most $b$ boundary vertices and at most $s$ interior vertices, then
$w_v=O(b+s)$, and the cost is $O\bigl(|V|(b+s)^2\bigr)$, linear in $|V|$.
\end{theorem}

\begin{proof}
When an interior vertex is eliminated, its row and column involve only vertices of its own snarl and
that snarl's boundary. Earlier snarls have been reduced to Schur complements on boundaries shared with
it. Proposition~\ref{ch16:prop:schur} and Lemma~\ref{ch16:lem:piv} give correctness.
\end{proof}

\paragraph{Arithmetic.} Over $\mathbb Q$ the entries of the Schur complements are rationals whose size
grows along the chain. For exact integers on large graphs we eliminate modulo several primes and
recombine. A zero pivot modulo a prime has probability about $|V|/p$, and it is detected. Two
independent orders, snarl order and minimum degree, give two independent checks.

\section{Computations}\label{ch16:sec:comp}

All numbers are produced by \texttt{bio14\_verify.py}, which writes \texttt{bio14\_output.txt}. The
battery runs in $44$ seconds.

\paragraph{The identities.} Proposition~\ref{ch16:prop:schur} is checked on $60$ random Eulerian digraphs
with random interiors, including the sign conditions on $S$. Proposition~\ref{ch16:prop:forest} is checked on
$40$ random interiors against exhaustive enumeration of out-arc choices. The BEST theorem is checked on
$40$ random digraphs with up to $12$ arcs against exhaustive enumeration of Eulerian circuits.

\paragraph{Synthetic genome graphs.} A chain of snarls between boundary vertices $b_i$ is traversed by
$2$--$4$ haplotypes, as one closed walk. Each snarl has an interior pool of four vertices and a nested
bubble with its own pool. The snarl order eliminates the nested interior, the nested boundary, the
snarl interior, and then the entry vertex. Table~\ref{ch16:tab:syn} compares it with dense Bareiss
elimination and with the forced-arc contraction of our spectral-fibres paper Chapter~\ref{chap:ch12}. All three are
exact and agree.

\begin{table}[t]
\centering\small
\begin{tabular}{rrrrrl}
\toprule
snarls & haplotypes & vertices & arcs & max front & $t_r(G)$\\
\midrule
10 & 2 & 74 & 126 & 4 & $10$ digits; three methods agree\\
20 & 3 & 162 & 360 & 6 & $37$ digits; three methods agree\\
40 & 3 & 311 & 691 & 7 & $79$ digits; three methods agree\\
60 & 4 & 516 & 1\,365 & 8 & $159$ digits; snarl order $=$ contraction\\
\bottomrule
\end{tabular}
\caption{Exact arborescence counts on synthetic genome graphs. The three methods are snarl-order
elimination, dense Bareiss (skipped above $320$ vertices) and forced-arc contraction.}\label{ch16:tab:syn}
\end{table}

\paragraph{Scaling.} Table~\ref{ch16:tab:scale} runs the snarl order and the minimum-degree order modulo two
primes on graphs of up to $39\,714$ vertices. They agree modulo both primes. The front stays at $8$ in
snarl order and at $6$--$7$ in minimum-degree order, and the time per vertex is essentially constant.

\begin{table}[t]
\centering\small
\begin{tabular}{rrrrrrc}
\toprule
snarls & vertices & arcs & front (snarl) & front (min-deg.) & $\mu$s / vertex & agree mod $p_1,p_2$\\
\midrule
500 & 3\,981 & 8\,638 & 8 & 6 & 9.2 & yes\\
1\,500 & 11\,949 & 26\,197 & 8 & 6 & 9.3 & yes\\
5\,000 & 39\,714 & 86\,967 & 8 & 7 & 10.7 & yes\\
\bottomrule
\end{tabular}
\caption{Linear time on synthetic genome graphs ($3$ haplotypes), $p_1=2^{31}-1$, $p_2=2\,147\,483\,629$.}
\label{ch16:tab:scale}
\end{table}

\paragraph{$\phi$X174.} The de Bruijn graph of the circular genome at word length $k$ has one vertex per
distinct $k$-mer and $5\,386$ arcs. It has no snarl tree given in advance, so we use minimum-degree order.
Table~\ref{ch16:tab:phix} compares it with forced-arc contraction; the two agree exactly. The front measures
how tangled the graph is. At $k=7$, where most $7$-mers recur by chance, it is $196$, and the elimination
takes $35$ seconds in exact rationals. At $k=12$ it is $2$, and $t_r=2$.

\begin{table}[t]
\centering\small
\begin{tabular}{rrrrr}
\toprule
$k$ & vertices & arcs & max front & $\log_{10}t_r(G)$\\
\midrule
7 & 4\,115 & 5\,386 & 196 & 336.5\\
8 & 4\,928 & 5\,386 & 80 & 128.2\\
10 & 5\,339 & 5\,386 & 12 & 12.6\\
12 & 5\,384 & 5\,386 & 2 & 0.3\\
\bottomrule
\end{tabular}
\caption{Exact arborescence counts of the single-stranded de Bruijn graphs of $\phi$X174, by
minimum-degree Schur elimination and by forced-arc contraction.}\label{ch16:tab:phix}
\end{table}

\section{What this does not do}\label{ch16:sec:not}

\begin{itemize}
\item \emph{It does not touch the double-stranded count.} $\#\DS$ is a sum of determinants over the points
of a lattice in a box, and it is $\sP$-hard Chapter~\ref{chap:ch10}. A Schur complement is a function of one
matrix. It does not commute with a sum over lattice points, which is why Chapter~\ref{chap:ch09}'s frontier had to carry
combinatorial states instead of boundary matrices Chapter~\ref{chap:ch09}.
\item \emph{It does not change the complexity class of the single-stranded count}, which BEST already
places in polynomial time. It turns a cubic dense computation into a linear sparse one when the snarl
width is bounded. The $\phi$X174 graphs at small $k$ show that the width need not be small.
\item \emph{No truncation.} Singular-value truncation of bond dimensions, as proposed, would make the
count approximate. That defeats the purpose of an exact count, and it is unnecessary: the exact bond
dimension is the boundary size.
\item \emph{Not addressed:} phasing constraints for diploid and polyploid genomes, and measurements on
the human MHC or on centromeric repeats, as proposed. We have no graphs for those loci, and phasing is a
different problem from counting.
\end{itemize}

\section*{Ancillary file}
\texttt{bio14\_verify.py} (standard library; imports \texttt{ds\_fermion.py} of Chapter~\ref{chap:ch09} and
\texttt{spectral\_fibres.py} of Chapter~\ref{chap:ch12}) runs checks (M), (F), (B), (E), (P) and (S). It writes
\texttt{bio14\_output.txt}.
